\clearpage

\section{SMC con atenuación de \textit{chattering}}

Considérese el sistema de segundo orden
\[
  \dot x_1=x_2,\qquad
  \dot x_2=f(x,t)+u,\qquad
  y=x_1.
\]

Para el error de seguimiento
\[
  e=r-y=r-x_1,
\]
se tiene
\[
  \dot e=\dot r-x_2,
  \qquad
  \ddot e=\ddot r-f(x,t)-u.
\]

Se define la superficie de deslizamiento integral
\[
  \boxed{\sigma=\dot e+m_1e+m_2\int_0^t e(\lambda)\,d\lambda}
  \qquad m_1>0,\quad m_2>0
\]

Su derivada es
\begin{align*}
  \dot\sigma
  &=\ddot e+m_1\dot e+m_2e\\
  &=\ddot r+m_1\dot r-f(x,t)-m_1x_2+m_2e-u.
\end{align*}

Suponiendo que existe un número \(L>0\) tal que
\[
  \left|\ddot r+m_1\dot r-f(x,t)\right|\leq L,
\]
se define
\[
  \boxed{V_{\mathrm{eq}}=L+m_1|x_2|+m_2|e|}
\]

La función de Lyapunov
\[
  V=\frac{1}{2}\sigma^2
\]
tiene como derivada
\[
  \dot V
  \leq
  V_{\mathrm{eq}}|\sigma|-\sigma u.
\]

Se propone la ley de control
\[
  \boxed{u=K V_{\mathrm{eq}}\operatorname{sign}(\sigma)},
  \qquad K>1.
\]
Entonces,
\[
  \dot V
  \leq
  -(K-1)V_{\mathrm{eq}}|\sigma|,
\]
por lo que la superficie es atractiva bajo las cotas asumidas.

\subsection{Aplicación al modelo PVTOL}

Para cada canal se definen los errores y sus integrales:
\begin{align*}
  e_x&=r_x-x, &
  I_x&=\int_0^t e_x(\lambda)\,d\lambda,\\
  e_z&=r_z-z, &
  I_z&=\int_0^t e_z(\lambda)\,d\lambda,\\
  e_\theta&=r_\theta-\theta, &
  I_\theta&=\int_0^t e_\theta(\lambda)\,d\lambda.
\end{align*}

Las superficies de deslizamiento son
\begin{align*}
  \sigma_x&=\dot e_x+m_{x1}e_x+m_{x2}I_x,\\
  \sigma_z&=\dot e_z+m_{z1}e_z+m_{z2}I_z,\\
  \sigma_\theta&=\dot e_\theta
  +m_{\theta1}e_\theta+m_{\theta2}I_\theta.
\end{align*}

Como
\[
  \dot e_x=\dot r_x-\dot x,\qquad
  \dot e_z=\dot r_z-\dot z,\qquad
  \dot e_\theta=\dot r_\theta-\dot\theta,
\]
las superficies también pueden escribirse como
\begin{align*}
  \sigma_x
  &=\dot r_x-\dot x+m_{x1}(r_x-x)+m_{x2}I_x,\\
  \sigma_z
  &=\dot r_z-\dot z+m_{z1}(r_z-z)+m_{z2}I_z,\\
  \sigma_\theta
  &=\dot r_\theta-\dot\theta
  +m_{\theta1}(r_\theta-\theta)+m_{\theta2}I_\theta.
\end{align*}

Para cada canal se considera
\begin{align*}
  V_{\mathrm{eq},x}
  &=L_x+m_{x1}|\dot x|+m_{x2}|e_x|,\\
  V_{\mathrm{eq},z}
  &=L_z+m_{z1}|\dot z|+m_{z2}|e_z|,\\
  V_{\mathrm{eq},\theta}
  &=L_\theta+m_{\theta1}|\dot\theta|
  +m_{\theta2}|e_\theta|.
\end{align*}

Las entradas de control son
\begin{align*}
  v_x&=K_x V_{\mathrm{eq},x}
  \operatorname{sign}(\sigma_x),\\
  v_z&=K_z V_{\mathrm{eq},z}
  \operatorname{sign}(\sigma_z),\\
  \tau&=K_\theta V_{\mathrm{eq},\theta}
  \operatorname{sign}(\sigma_\theta).
\end{align*}

Los parámetros deben satisfacer
\[
  m_{x1}>0\qquad m_{x2}>0\qquad
  m_{z1}>0\qquad m_{z2}>0\qquad
  m_{\theta1}>0\qquad m_{\theta2}>0,
\]
\[
  K_x>1\qquad K_z>1\qquad K_\theta>1.
\]